% RESEARCH_DRAFT_SHA256: 187a278b307d80785b74ec3514bdb799e982d969b05240033b647f55e68c1b99
% Layout and source-reference wrappers expanded for the web renderer.
\documentclass[11pt]{article}
\usepackage[margin=25mm]{geometry}
\usepackage{fontspec,unicode-math}
\setmainfont{Latin Modern Roman}
\setmathfont{Latin Modern Math}
\usepackage{amsmath,amssymb,mathtools,xurl,hyperref}
\hypersetup{colorlinks=true,urlcolor=blue,linkcolor=blue}
\setlength{\parindent}{0pt}
\setlength{\parskip}{5pt}
\setlength{\emergencystretch}{3em}
\newcommand{\sref}[2]{\hyperlink{source:#1:#2}{[S#2]}}
\newcommand{\cataloguescope}[1]{\paragraph{Source question.}#1}
\newcommand{\floor}[1]{\left\lfloor#1\right\rfloor}
\begin{document}

\section*{Bounded subdeterminants: rational vertices and two exponential enumeration traps}
\noindent UnsolvedMath ID 30004899; OWR-8415356-004. Research attempt dated 11 October 2026.
\paragraph{Result in scope.}This manuscript contains a source-specific attempt and complete proofs of the stated auxiliary calculations. It makes no new solution or novelty claim for the entire target.
\section{Original problem and preserved definition}
The following definition and source list are reproduced from the repository. Any earlier inspection dates and reported literature claims in this preserved text describe that source record, not additional verification by this attempt.
% BEGIN PRESERVED DEFINITION
\subsection{Definitions and mathematical statement}

\paragraph{Shared notation.}Standard number systems are denoted by $\mathbb N,\mathbb Z,\mathbb Q,\mathbb R,\mathbb C$; index ranges, computational models and specialized notation follow the cited problem source and the definitions below.


An integer program optimizes a linear objective over $\{x\in\mathbb Z^n:Ax\leq b\}$ with finitely encoded integer constraint data. A square subdeterminant is the determinant of any square submatrix of $A$. The parameter $\Delta$ bounds their absolute values and is fixed independently of the dimensions. Polynomial time means polynomial in the total binary input length, including $b$ and the objective.
\paragraph{Statement recorded in the supplied source.}
For a fixed integer $\Delta\geq3$, can integer programs whose constraint matrices have every square subdeterminant bounded in absolute value by $\Delta$ be solved in polynomial time? \href{https://doi.org/10.4171/owr/2021/53}{[S2]}
\paragraph{Source status.}
The supplied catalogue labels this record partially solved. Its source-status metadata was checked against the supplied snapshot on 2026-10-11. This import does not independently establish current open status or mathematical correctness. The catalogue literature assessment is dated 2026-08-22.
\subsection{Short English statement}
Determine whether a constant bound on every constraint subdeterminant makes integer programming polynomial-time solvable.
\subsection{Sources}
\begin{itemize}
\item[S1] Ulam AI and the UnsolvedMath Contributors, UnsolvedMath: A Curated Collection of Open Mathematics Problems, record 30004899 (OWR-8415356-004). Catalogue curation and metadata: CC BY 4.0. The supplied source question is retained; definitions and the English summary are editorial additions.
\url{https://huggingface.co/datasets/ulamai/UnsolvedMath}.
\item[S2] Combinatorial Optimization (2022), Oberwolfach Reports, DOI 10.4171/owr/2021/53, pp. 2893--2954
\url{https://doi.org/10.4171/owr/2021/53}.
\item[S3] Combinatorial Optimization, Oberwolfach Reports 53/2021, primary report PDF
\url{https://ems.press/content/serial-article-files/46931}.
\textit{Underlying Oberwolfach Reports material retains its source-specific licence, including CC BY-SA 4.0 where applicable; see THIRD\_PARTY\_NOTICES.md.}
\end{itemize}
% END PRESERVED DEFINITION
\subsection{Verbatim supplied catalogue question}
For a fixed integer $\Delta\geq3$, can integer programs whose constraint matrices have every square subdeterminant bounded in absolute value by $\Delta$ be solved in polynomial time?

\section{Attempted finite-denominator algorithm}
Every square minor is bounded by a fixed $\Delta$, so in particular
each entry of $A$ is bounded. It is tempting to enumerate the
few possible fractional vertex types or search a bounded integer box.
We derive the genuine denominator statement and then show why neither
enumeration argument has a polynomial input-length bound.

\paragraph{Vertex denominators are bounded.}
At a vertex $x$ of a full-dimensional polyhedron $Ax\le b$,
the active row vectors span $\mathbb R^n$. Otherwise choose
nonzero $h$ orthogonal to all active rows. Because there are finitely
many inactive inequalities, both $x+th$ and $x-th$ remain
feasible for sufficiently small $t>0$: choose $t$ below every
positive slack divided by the corresponding nonzero directional
absolute value. Then $x$ is their midpoint and is not a vertex,
a contradiction. Choose $n$ independent active rows, forming an
integer nonsingular matrix $B$. The equality $Bx=b_B$ gives
$x=B^{-1}b_B=\operatorname{adj}(B)b_B/\det B$.
The adjugate entries and $b_B$ are integral. Thus each coordinate
denominator divides the nonzero integer $\det B$, whose absolute
value is at most $\Delta$. In particular all vertices lie in
$D^{-1}\mathbb Z^n$ for $D=\operatorname{lcm}(1,\ldots,\Delta)$.
This does not imply a bounded number of vertices or nearby lattice
points, or that coordinatewise rounding preserves inequalities.

\paragraph{Fractional-vertex rounding fails already at determinant three.}
Minimize $x+y$ subject to
\[
2x+y\ge1,\qquad x+2y\ge1,\qquad0\le x,y\le1.
\]
The linear optimum is $(1/3,1/3)$ with value $2/3$:
adding the first two inequalities gives $x+y\ge2/3$,
and that point attains it. The constraint rows are
$(-2,-1),(-1,-2),(1,0),(0,1),(-1,0),(0,-1)$.
All one-by-one minors have absolute value at most two; the
two special rows have determinant three, a special row and a
coordinate row give absolute determinant at most two, and two
coordinate rows give at most one. Hence every square minor is
bounded by three. Rounding both coordinates down gives $(0,0)$,
which violates both covering constraints. Rounding both up gives
$(1,1)$ of value two, whereas $(1,0)$ and $(0,1)$ are
feasible integer points of value one and no integer value below
one is feasible. Even this small denominator does not make either
uniform rounding rule an exact integer optimizer.

\paragraph{Enumerating all vertices or all box points is not enough.}
The cube $[0,1]^n$ has exactly $2^n$ vertices: a point with
a coordinate strictly between zero and one is a midpoint of a
small coordinate perturbation, whereas a zero-one point cannot be
a nontrivial convex combination of other cube points, coordinate
by coordinate. Its matrix consists of rows $\pm e_i$ and
all minors belong to $\{0,1,-1\}$: a square submatrix with
a zero row or repeated nonzero column position has determinant
zero; otherwise it is a signed permutation matrix. Thus even
the stronger determinant-one promise permits exponentially many
vertices. Listing all LP bases or vertices is not polynomial.

Likewise the one-dimensional interval $0\le x\le B$ has a
determinant-one constraint matrix and $B+1$ integer points. If
$B$ is binary encoded, listing them is exponential in $\log B$.
These simple instances are themselves easy to optimize; they disprove
only the proposed enumeration bounds, not polynomial solvability of
the promised class. A residue-class list of size $D^n$ also has
no polynomial bound at fixed $D>1$. We do not claim that all
those residue patterns can independently occur under the full minor
promise; using this loose list already loses the desired running time.

\section{Audit}
The primary bounded-subdeterminant discussion on printed page 2912
was inspected. It treats two nonzeros per row as an additional
hypothesis, which is not silently imposed here. The audit checked
all minor sizes in the two-variable example, full-dimensionality,
binary input length and the distinction between determinant bounds
and coordinate bounds. No integer-programming solver or vertex
enumerator was run. An algorithm exploiting the global restrictions
among all minors, rather than the displayed loose enumerations,
remains to be supplied.

\section{Process and contributions}
Mathematics and model-assisted adversarial self-audit: OpenAI gpt-6-astra, ultra reasoning, as dispatched for this research batch; exact serving revision not exposed. Codex (GPT-6; exact serving revision not exposed) materially assisted with orchestration and format review. No human mathematical authorship, mathematical edits, or independent human proof checking is asserted.
Contribution prompt A was applied to the full definition and supplied source row, then prompt B was applied adversarially to the derived results. The item-specific audit above records the scope checks and rejected extrapolations. Local assembly and structural checks do not establish mathematical correctness. The source catalogue status is preserved in TOP\_PROBLEM and is separate from this attempt status.
\paragraph{Integration review.} Codex (GPT-6; exact serving revision not exposed) reviewed the written arguments and expanded layout and source-reference wrappers. Repository metadata and the production web renderer were checked using Python 3.13 and JavaScript with MathJax 3.2.2. This model-assisted review is not independent human proof verification; no PDF compilation or formal proof checking is asserted.

\section{Completion Estimate}
\noindent\textbf{COMPLETION: 0\%}
The proved material is a rederivation of known or auxiliary facts rather than a substantiated new advance on the original target. Bounded vertex denominators do not justify rounding or polynomial enumeration; no general fixed-Delta integer-programming algorithm is constructed.
\section{References}
The complete original Sources subsection above is retained, including catalogue attribution and licence notices. The additional list records primary sources actually consulted for this attempt on 11 October 2026; uncited source links reproduced above were not all rechecked.
\begin{itemize}
\item Oberwolfach Report 53/2021, Combinatorial Optimization, original question and surrounding discussion read from the primary PDF text on 11 October 2026. Historical statements are not a current-status survey. \url{https://ems.press/content/serial-article-files/46931}.
\end{itemize}
\section{Final status}
\textbf{Unresolved in this attempt.} Bounded vertex denominators do not justify rounding or polynomial enumeration; no general fixed-Delta integer-programming algorithm is constructed.
\end{document}
